% ============================================================
%  MÓDULO 13 — ATIVAÇÃO DE AGENTES E SUBAGENTES (minuciosa)
%  Pipeline determinístico único + ativação por categoria.
%  Teoria, função, cálculo e gates, com citações DOI verificadas
%  (cronologia 1..7 confere com a biblioteca do Módulo 11).
% ============================================================

\chapter{Ativação de Agentes e Subagentes}
\label{mod13}

\roteirodidatico
{Uma descrição de agente no catálogo não significa que ele tenha sido chamado. Este módulo distingue o registro disponível das etapas que conduzem uma solicitação a uma execução.}
{\emph{Gatilho} inicia o fluxo; \emph{roteamento} escolhe a rota; \emph{permissão} autoriza uma operação; \emph{feedback} devolve resultado; \emph{registro} conserva o rastro.}
{Acompanhe os sete estágios com uma tarefa de exemplo e identifique, em cada um, a entrada, a decisão, a saída e o ponto em que a execução pode parar.}
{Modelos de pontuação e cálculos ilustrativos explicam regras, mas não são medidas empíricas de precisão sem uma amostra, referência de comparação e avaliação reportada.}
{Que evidência distingue um agente apenas catalogado de um agente selecionado, chamado e concluído com sucesso?}

\casoguiado{estados distintos de ativação}
{Em uma sequência ilustrativa, a ficha existe no catálogo; a configuração a registra; o aplicativo a carrega; uma tarefa a seleciona; o executor retorna uma saída. Cada etapa precisa de uma evidência própria. Se só houver a ficha, a conclusão correta é ``documentado''. Se houver um log de chamada, pode-se dizer ``invocado naquela execução'', ainda sem afirmar que a resposta foi correta.}

\begin{resultado}
ESCOPO. Este módulo apresenta um modelo de referência para analisar possíveis
rotas de ativação, desde um pedido até a execução e o registro. A sequência
(1) gatilho, (2) roteamento, (3) correspondência de capacidades, (4) verificação
de permissões, (5) execução, (6) retorno e (7) registro é uma decomposição
didática: os caminhos reais podem variar entre aplicativo, configuração, skill,
integração e ambiente. Não se deve presumir que todos os sete passos ocorram
sempre nem que a seleção seja determinística. As seções distinguem o fluxo
documentado das partes que precisam de confirmação operacional e dos exemplos
matemáticos, que são ilustrativos até serem avaliados com execuções observadas.
categoria.
\end{resultado}

% ============================================================
\section{Sistema geral de ativação — o pipeline de 7 etapas}
% ============================================================

\begin{processo}
ETAPA 1 — GATILHO. No modelo de análise deste capítulo, o gatilho é o evento que
leva uma tarefa a ser considerada para execução: por exemplo, uma skill
solicitada, um perfil nomeado, um comando ou uma tarefa publicada no Blackboard.
O fluxo concreto pode iniciar de outra forma; identifique a interface e os dados
de entrada no caso estudado antes de chamar o evento de gatilho.
\end{processo}

O gatilho só é legítimo se nomear o que deve ser feito; o vocabulário de
cada card (descrição, skills, tags) foi construído para casar com esses
gatilhos. Isso é deliberado: o conceito de agente neste ecossistema segue a
tríade que Wooldridge e Jennings tornaram canônica — reatividade à mudança do
ambiente, proatividade orientada a objetivos e habilidade social de
interação \citref{WOOLDRIDGE, M.; JENNINGS, N. R. Intelligent agents: theory and practice. \emph{The Knowledge Engineering Review}, v.~10, n.~2, p.~115-152, 1995}{10.1017/S0269888900008122}{é o lastro teórico da especificação do gatilho: um agente só é um agente se é reativo, proativo e social — propriedades que o pipeline de ativação transforma em regras de roteamento (reagir à tarefa publicada, perseguir o objetivo declarado e interagir via Blackboard/MetaBus)}{The concept of an agent has become important in both artificial intelligence and mainstream computer science}{O conceito de agente tornou-se importante tanto na inteligência artificial quanto na ciência da computação.}{WOOLDRIDGE; JENNINGS, 1995, p.~115--152}.\begin{processo}
ETAPA 2 — ROTEAMENTO. O orquestrador decide o caminho de ativação em ordem de
preferência rígida: (a) skill marcada como API-alvo da tarefa; (b) agente
nomeado explicitamente (rota inline em \pth{opencode.json}); (c) subagente do
catálogo cujas capacidades casem com o gatilho; (d) fallback: ``executante
genérico'' sem especialização. A escolha (b) roda direto; a (c) entra no
Blackboard para voluntariado.
\end{processo}

\begin{processo}
ETAPA 3 — CASAMENTO DE CAPACIDADES. Para a rota (c), calcula-se o escore de
ativação de cada card voluntário. Seja $K$ o conjunto de tokens-chave do
gatilho (remoção de stopwords) e $C_a$ as capacidades declaradas do agente
$a$; o escore é a interseção em união (Jaccard ponderado):
$A(a) = |K \cap C_a| / |K \cup C_a|$, $0 \le A \le 1$. Empates são
desfeitos por prioridade de categoria (curation \textgreater research
\textgreater audit) e, em último caso, por ordem alfabética — nunca por
aleatoriedade, para preservar a reprodutibilidade.
\end{processo}

\begin{processo}
ETAPA 4 — GATE DE PERMISSÃO. Antes da execução, o card é validado pelo
bootstraper (SPEC-935-R608: frontmatter tolerante, descrição não vazia,
capacidades não vazias) e as permissões declaradas no card são aplicadas.
Card ilegível não executa: erro visível, não silencioso (fail-closed).
\end{processo}

\begin{processo}
ETAPA 5 — EXECUÇÃO. O agente recebe o gatilho + contexto do MetaBus e executa
na rota escolhida (inline, skill, MCP). Se o agente estiver marcado com
política de invocação implícita proibida, o orquestrador lê o
\pth{SKILL.md} correspondente e executa suas instruções no contexto atual,
sem re-disparo de habilidade.
\end{processo}

A distribuição da tarefa entre múltiplos especialistas voluntários — e não a
execução por um único programa fixo — é o coração arquitetural do Core. A
justificativa teórica vem da maturidade do campo de sistemas multiagente: a
coordenação por delegação e o reuso de capacidades declaradas são os
princípios que o roteiro de pesquisa do campo consolida, e é esse reuso que o
Escore de Ativação mede \citref{JENNINGS, N. R.; SYCARA, K.; WOOLDRIDGE, M. A roadmap of agent research and development. \emph{Autonomous Agents and Multi-Agent Systems}, v.~1, n.~1, p.~7-38, 1998}{10.1023/A:1010090405266}{é a moldura em que o voluntariado por capacidades se apoia: o campo consolida que agentes autônomos são construídos em torno de capacidades reutilizáveis e interação estruturada — exatamente o que o casamento de capacidades e o MetaBus implementam}{It aims to identify key concepts and applications, and to indicate how they relate to one-another}{O objetivo é identificar conceitos-chave e aplicações e indicar como eles se relacionam entre si.}{JENNINGS, 1998, p.~7--38}.\begin{processo}
ETAPA 6 — FEEDBACK. O resultado da execução volta ao orquestrador com o
vínculo especificação-teste (gate SDD/TDD): se a tarefa alterou código ou
manual, o ciclo de evolução é registrado com score e lições. O resultado não
é ``verificado'' sem validação externa.
\end{processo}

\begin{processo}
ETAPA 7 — REGISTRO. Cada ativação que mudou estado gera um ciclo em
\pth{evolution/cycles.json} — objetivo, mudanças, score, lições e
round\_id sem colisão com specs (SPEC-935-R611). Reprogramar a execução
exige fonte de verdade: o mesmo gatilho + mesmo catálogo produz a mesma
rotação (determinismo).
\end{processo}

O requisito de determinarismo — mesmo gatilho, mesmo catálogo, mesmo manual —
é o padrão mínimo de reprodutibilidade para alegações computacionais. Como o
manual é gerado a partir do catálogo, a reprodutibilidade da ativação é
herdada pela reprodutibilidade da geração; quando a replicação integral é
impossível, a reprodutibilidade serve como padrão mínimo de julgamento \citref{PENG, R. D. Reproducible research in computational science. \emph{Science}, v.~334, n.~6060, p.~1226-1227, 2011}{10.1126/science.1213847}{ancora a política de determinismo do pipeline de ativação e do gerador do manual: reprodutibilidade é o padrão mínimo para julgar alegações computacionais quando a replicação completa não é viável — princípio que o livro aplica ao próprio livro}{Reproducibility has the potential to serve as a minimum standard for judging scientific claims when full independent replication of a study is not possible}{A reprodutibilidade tem o potencial de servir como padrão mínimo para julgar alegações científicas quando a replicação independente completa de um estudo não é possível.}{PENG, 2011, p.~1226--1227}.\begin{arquitetura}
FUNÇÕES ÚNICAS POR ETAPA. E1 gatilho: `parsea gatilho e publica tarefa` em
\pth{mci/mcp_server.py} (mci\_post\_task). E2 roteamento: `decide caminho`
em \pth{marceloclaro/orchestrator.py}. E3 casamento: `escore de ativação` em
\pth{mci/blackboard.py} (matching de capacidades). E4 gate:
`parse/card/permissão` em \pth{mci/agent_registry_bootstrap.py}. E5 execução:
`executa skill/subagente` (rota MCP/inline). E6 feedback: `valida gate
SDD/TDD`. E7 registro: `record()` em \pth{evolution/cycles.py}
(round\_id único, R611).
\end{arquitetura}

\begin{aviso}
AVISO/ANTI-OVERCLAIM. O pipeline descreve COMO a ativação acontece; descrever
o mecanismo não garante a qualidade do resultado. Um agente bem ativado pode
produzir má resposta; a diferença é auditável pela trilha E1--E7, não por
afirmação. (R110.)
\end{aviso}

% ============================================================
\section{Ativação por categoria — lote por categoria}
% ============================================================

\begin{descricao}
ORGANIZAÇÃO. As seções a seguir aplicam o pipeline de 7 etapas a CADA
categoria do catálogo (mod-10). Para cada categoria: gatilho típico, rota
dominante, função do agente, cálculo específico, justificativa teórica,
impactos e limites e aviso. As categorias são as 12+ do gerador
\pth{gen\_mod10.py} — orquestração, acadêmica, pesquisa, auditoria, jurídica,
dados, inferência, engenharia, testes, literária e cura (paisagem; mapeia o
card do agente-alvo e recomenda capacidades, sem substituir o pipeline E1--E7),
integração, MIRA, MASWOS e utilidade.
\end{descricao}

\subsection{Orquestração}
\elemento{ativ-orquestracao}{\metainfo{Categoria}{orquestracao} \hfill \metainfo{Rota}{marceloclaro/orchestrator.py}}

\begin{descricao}
GATILHO. Qualquer pedido de múltiplos passos; rota (b) inline. FUNÇÃO:
coordenar, não executar. CÁLCULO: também executa o E3 para escolher o
subagente com maior escore $A(a)$; mantém a trilha E1--E7 como registro.
JUSTIFICATIVA: reatividade/proatividade/socialidade (Wooldridge; Jennings
1995) + oportunismo blackboard (Nii 1986). LIMITE: o orquestrador depende da
qualidade dos cards — daí os gates R608/R609 que verificam a leitura e a
normalização das fichas; um card legível ainda não comprova a capacidade de
execução do agente.
\end{descricao}

\subsection{Acadêmica (artigos/TCC/dissertação)}
\elemento{ativ-academica}{\metainfo{Categoria}{academic} \hfill \metainfo{Rota}{skills/pesquisa-artigo-qualis-a1}}

\begin{descricao}
GATILHO: ``artigo'', ``TCC'', ``dissertação'', ``banca''. ROTA: skill
curada + subagentes 00--44. FUNÇÃO: cumprir o ciclo SDD/TDD editorial
(escopo, busca, evidências, estrutura, revisão, QA). CÁLCULO: a cobertura de
cada etapa é $1 - k/n$ (k critérios abertos em n criteriosos da rubrica
MASWOS). JUSTIFICATIVA: validação externa como única fonte de mérito
(Ioannidis 2005; R110). LIMITE: não promete aprovação nem Qualis.
\end{descricao}

\subsection{Pesquisa (literatura/evidências)}
\elemento{ativ-pesquisa}{\metainfo{Categoria}{research} \hfill \metainfo{Rota}{skills/pesquisador-universal}}

\begin{descricao}
GATILHO: ``revisão sistemática'', ``buscar fontes'', ``meta-análise''.
FUNÇÃO: busca e curadoria com download open access e verificação de fontes.
CÁLCULO: grau de cobertura da living review = fontes verificadas por PMID/DOI
dividido pelo total de candidatas. JUSTIFICATIVA: reprodutibilidade mínima
(Peng 2011) aplicada a evidências. LIMITE: corpus é direcionado pela query;
não substitui consenso de especialistas.
\end{descricao}

\subsection{Auditoria e qualidade}
\elemento{ativ-auditoria}{\metainfo{Categoria}{audit} \hfill \metainfo{Rota}{agente auditor + honest-critic}}

\begin{descricao}
GATILHO: ``audite'', ``confere'', ``revisão de qualidade''. FUNÇÃO: separar
cobertura de mérito; aplicar fail-closed. CÁLCULO: colunas da auditoria —
alegação, evidência, validação; ausência de validação externa gera
``inconclusivo/requer\_escalonamento''. JUSTIFICATIVA: provável falsidade de
alegações não validadas (Ioannidis 2005). LIMITE: auditor não conhece o
mundo real.
\end{descricao}

\subsection{Jurídica (auxjuris)}
\elemento{ativ-juridica}{\metainfo{Categoria}{legal} \hfill \metainfo{Rota}{agents/auxjuris-*}}

\begin{descricao}
GATILHO: ``cláusula'', ``petição'', ``minuta de ofício'', ``jurisprudência''.
ROTA: subagentes auxjuris (research, document\_summarizer, email\_drafter,
legal\_assistant). FUNÇÃO: apoiar pesquisa/digestão documental e redação de
minutas com revisão humana obrigatória. CÁLCULO: checklist de cobertura por
apontamento; texto sem revisão humana é classificado `rascunho`. LIMITE: não
emite parecer autônomo nem responde à diligência por conta própria.
\end{descricao}

\subsection{Dados e proveniência}
\elemento{ativ-dados}{\metainfo{Categoria}{data} \hfill \metainfo{Rota}{skills/engenharia de dados + FAIR}}

\begin{descricao}
GATILHO: ``datasets'', ``proveniência'', ``coleta''. FUNÇÃO: garantir que
todo dado tenha origem, versão, papel analítico e restrição documentados.
CÁLCULO: cobertura FAIR por dataset = proporção dos princípios F/A/I/R
atendidos com evidência. JUSTIFICATIVA: proveniência detalhada (Wilkinson
2016, R1.2). LIMITE: dados coletados de API pública podem quebrar sem aviso.
\end{descricao}

\subsection{Inferência estatística e computacional}
\elemento{ativ-inferencia}{\metainfo{Categoria}{inference} \hfill \metainfo{Rota}{agentes 20--28}}

\begin{descricao}
GATILHO: ``estatística'', ``modelagem'', ``ML'', ``benchmarking''.
FUNÇÃO: adequação de testes, modelagem formal e robustez de resultados.
CÁLCULO: baseline+ablação; generalização aferida por erro fora da amostra.
JUSTIFICATIVA: não declarar efeito sem controle (Ioannidis 2005; Peng 2011).
LIMITE: resultado estatístico é contingente ao desenho declarado.
\end{descricao}

\subsection{Engenharia e infraestrutura}
\elemento{ativ-engenharia}{\metainfo{Categoria}{engineering} \hfill \metainfo{Rota}{skills/cloud-* + coder}}

\begin{descricao}
GATILHO: ``construa'', ``implemente'', ``infra'', ``cloud''. FUNÇÃO:
codificar com protocolo SDD/TDD e operar componentes de infra.
CÁLCULO: gate = testes verdes / specs com teste (R610). JUSTIFICATIVA:
evidência executável antes de ``concluído'' (Ioannidis 2005). LIMITE:
ambiente mínimo declarado; sem commitar sem autorização.
\end{descricao}

\subsection{Testes}
\elemento{ativ-testes}{\metainfo{Categoria}{testing} \hfill \metainfo{Rota}{tests/test\_r*.py}}

\begin{descricao}
GATILHO: ``teste'', ``regressão'', ``pytest''. FUNÇÃO: RED-verde do ciclo
SDD/TDD e garantia de não-regressão (ex.: motor R608/R611).
CÁLCULO: cobertura de rotas = casos verdes / casos totais por família.
JUSTIFICATIVA: reprodutibilidade (Peng 2011) e fail-closed (Ioannidis 2005).
LIMITE: teste verde não generaliza para produção.
\end{descricao}

\subsection{Literária (projeto Molambudos)}
\elemento{ativ-literaria}{\metainfo{Categoria}{literary} \hfill \metainfo{Rota}{agentes literary-*}}

\begin{descricao}
GATILHO: ``trecho'', ``personagem'', ``voz autoral''. FUNÇÃO: aplicar
scanners literários e guardar voz/ética/tradução. CÁLCULO: score de
consistência de voz por marcador presente versus esperado. JUSTIFICATIVA:
curadoria auditável (Peng 2011); LIMITE: nenhum scanner aprova equivalência —
toda aprovação é humana.
\end{descricao}

\subsection{Cura da paisagem de agentes}
\elemento{ativ-curador}{\metainfo{Categoria}{curation-agent} \hfill \metainfo{Rota}{landscape-curator}}

\begin{descricao}
GATILHO: ``adicionar agente externo'', ``integrar nova coleção''. FUNÇÃO:
inventariar pontos de entrada externos e recomendar formato/capabilities do
card. CÁLCULO: paridade multilcoleção = cards próprios vs. manifests externos.
JUSTIFICATIVA: proveniência (Wilkinson 2016). LIMITE: card gerado é
candidato; ativação final exige o gate R608 (frontmatter válido).
\end{descricao}

\subsection{Integração (externos e MCP)}
\elemento{ativ-integracao}{\metainfo{Categoria}{integration} \hfill \metainfo{Rota}{integrations/*}}

\begin{descricao}
GATILHO: ``executor externo'', ``MCP'', ``orquestrar CLI''. FUNÇÃO: compor
CLIs externas (gemini, goose, plandex, reasonix, haystack) como executores
tolerantes — ausência do binário gera warn, nunca crash. CÁLCULO: saúde =
executores disponíveis / executores declarados. JUSTIFICATIVA: reuso de
capacidades (Jennings, Sycara, Wooldridge 1998). LIMITE: resultado externo
nunca é ``verificado'' sem validação do Core.
\end{descricao}

\subsection{MIRA (apresentações)}
\elemento{ativ-mira}{\metainfo{Categoria}{mira-agent} \hfill \metainfo{Rota}{agentes mira-*}}

\begin{descricao}
GATILHO: ``deck'', ``apresentação'', ``visual''. FUNÇÃO: pipeline de slides
(briefing, planejamento, animações, validação). CÁLCULO: conformidade de
deck = critérios do mira-validator atendidos / total. JUSTIFICATIVA:
determinismo de geração (Peng 2011). LIMITE: estética não mede eficácia de
comunicação.
\end{descricao}

\subsection{MASWOS (padrão nota 10)}
\elemento{ativ-maswos}{\metainfo{Categoria}{maswos-agent} \hfill \metainfo{Rota}{agentes 13--16, 44}}

\begin{descricao}
GATILHO: ``rubrica'', ``QA Qualis'', ``10/10''. FUNÇÃO: auditoria micro e
macro contra o padrão MASWOS. CÁLCULO: nota de rubrica = pontos atendidos /
pontos totais ponderados. JUSTIFICATIVA: separação cobertura/mérito
(Ioannidis 2005; R110). LIMITE: rubrica interna não é aceitação editorial.
\end{descricao}

\subsection{Utilidade (busca, git, docs, revisão)}
\elemento{ativ-utilidade}{\metainfo{Categoria}{utility} \hfill \metainfo{Rota}{git-manager, pypi-searcher, docs-writer, reviewer}}

\begin{descricao}
GATILHO: ``busque'', ``commite'', ``documente'', ``review''. FUNÇÃO:
serviços de suporte determinísticos. CÁLCULO: cobertura por serviço =
operações com rastro / operações totais. JUSTIFICATIVA: reprodutibilidade e
rastro (Peng 2011; Wilkinson 2016). LIMITE: convenções do repo mandam
(mensagem de commit, sem secret).
\end{descricao}

\begin{resultado}
RESUMO DO LOTE. Todas as categorias do catálogo foram percorridas com a
mesma chave de leitura (E1--E7 + função + cálculo + justificativa teórica +
limite + aviso). As citações usam apenas os DOIs verificados e ativos da
biblioteca do Módulo 11; nenhuma alegação de ``superioridade'' foi feita
(R110).
\end{resultado}

% ============================================================
\section{Como ativar na prática (exemplos executáveis)}
% ============================================================

\begin{metodo}
EXEMPLO 1 — ORQUESTRAR EXPANSÃO. Gatilho: ``expanda o módulo X com citações''.
E2 escolhe rota (b) orquestrador; E3 não se aplica a subagente nomeado; E4
gate de card do orquestrador; E6 exige teste/spec para cada mudança; E7 grava
ciclo com round\_id único (R611). O operador deve confirmar escopo antes de
gerar conteúdo em massa (núcleo primeiro, lote por categoria, etc).
\end{metodo}

\begin{metodo}
EXEMPLO 2 — ATIVAR AGENTE NOMEADO (``revisar com honest-critic''). E1 parsea
o nome; E2 rota (b) inline; E4 permissão lida do card; E5 execução; E6 o
resultado (parecer) alimenta o orquestrador; E7 registra se alterou manual.
A invocação de skill-filha com invocação implícita proibida exige que o
orquestrador leia o SKILL.md e execute no contexto — nunca re-disparo por
nome.
\end{metodo}

\begin{metodo}
EXEMPLO 3 — VOLUNTÁRIO (``preciso de um especialista em N''). E1 publica
tarefa; E3 calcula $A(a)$ para todos os cards; E5 os k = melhores
executam em paralelo; E6 combinados pelo orquestrador com nota de confiança do MetaBus;
E7 ciclo. Sem voluntário com $A>0$, cai no fallback (d) e o orquestrador
avisa o operador (transparência, não silêncio).
\end{metodo}

\begin{resultado}
APLICAÇÃO A TODOS OS MÓDULOS. Todo módulo do manual (1--13) que mencione um
agente/subagente herda este pipeline único: a ``ativação'' nunca é ambígua
entre módulos — é sempre E1--E7. Nos módulos 1--9 a referência é implícita;
a partir daqui (12--13) a referência é explícita e auditável contra a
biblioteca de DOIs verificados.
\end{resultado}

% ----------------------------------------------------------------------
%  N-13 — Leitura guiada do Módulo 13 (adicionada no ciclo R619)
% ----------------------------------------------------------------------
\section[Leitura guiada]{Leitura guiada — Módulo 13: ativar é percorrer sete etapas, não escolher um nome}

\elemento{Leitura guiada — Ativação de agentes}{\metainfo{ID}{N-13} \hfill \metainfo{Ciclo}{R619} \hfill \metainfo{Estilo}{ativar é percorrer sete etapas, não escolher um nome}}


\begin{descricao}
\textbf{O fenômeno.} Um nome no catálogo pode ser confundido com um executor
pronto para trabalhar. Na prática, há estados distintos: documentação presente,
configuração gerada, perfil carregado, rota selecionada, chamada realizada e
saída avaliada. O modelo de sete etapas abaixo serve para investigar esses
estados. A ordem e a presença de cada operação devem ser confirmadas no caminho
concreto da tarefa; este esquema não é uma garantia de que todo fluxo percorra
as mesmas etapas.
\end{descricao}

\begin{definicao}
\textbf{Grandezas e definições.}\par
(1) \emph{Gatilho}: pedido ou evento que pode iniciar a seleção de um caminho; não basta, por si só, para autorizar qualquer operação.\par
(2) \emph{Etapa}: uma das $7$ posições do pipeline (gatilho, roteamento, casamento, gate, execução, feedback, registro).\par
(3) \emph{Rota}: caminho de ativação; existem $4$ alternativas em ordem de preferência rígida.\par
(4) \emph{Ciclo de voluntariado}: publicar no Blackboard, aguardar candidatos e ponderar --- o custo da rota que passa pelo quadro.\par
(5) \emph{Voluntário}: agente cujas capacidades casam com o gatilho e se candidatam à tarefa postada.
\end{definicao}

\begin{metodo}
\textbf{Como analisar uma rota.} Para cada caminho encontrado no código,
registre: (1) qual evento o inicia; (2) que regra escolhe o executor; (3) que
permissões são consultadas; (4) que chamada efetivamente ocorre; (5) que saída
fica registrada. Se as regras e entradas forem fixas e não houver estado externo
variável, a decisão de roteamento pode ser determinística. Respostas de modelos,
serviços de rede, concorrência e estado compartilhado podem tornar o resultado
da execução não determinístico mesmo quando a seleção inicial se repete. Um
fallback só pode ser descrito como comportamento do sistema quando sua condição
e seu executor estiverem identificados na implementação.
\end{metodo}

\begin{resultado}
\textbf{Exemplo hipotético.} Suponha que, em um conjunto de tarefas, a rota
que usa o quadro seja escolhida em $p=0{,}2$ dos casos. Se ela acrescentar em
média $c=2$ operações observáveis em relação à rota direta, o custo médio
adicional estimado é $p\,c = 0{,}2\times2=0{,}4$ operação por tarefa. O cálculo
ilustra como estimar custo; os valores são pressupostos didáticos, não medições
do Core. Para estimar custo real, defina o que conta como operação, registre uma
amostra de execuções e reporte distribuição e falhas, não apenas a média.
\end{resultado}

\begin{resultado}
\textbf{Verificação.} A existência de uma etapa descrita no livro não demonstra
que ela esteja implementada. Para conferi-la, localize o ponto de entrada, siga
as chamadas e examine os testes ou registros aplicáveis. Um contador de
marcadores editoriais mede presença de texto, não cobertura do código nem
maturidade operacional.
\end{resultado}

\begin{aviso}
\textbf{Limite de validade.} Uma regra de roteamento simples pode facilitar
auditoria e, ainda assim, selecionar um executor inadequado. Para comparar
políticas de seleção seriam necessários casos representativos, uma referência
de resposta, métricas definidas antes do ensaio e repetição sob condições
controladas. Testes de código podem conferir regras implementadas; não demonstram
que a política produz melhores respostas. Até que essa avaliação exista, trate
qualquer vantagem de desempenho como hipótese.
\end{aviso}

\begin{aviso}
\textbf{Exercícios.}\par
(E1) Desenhe o caminho de uma tarefa que é ativada por um \emph{skill} (rota a) e identifique em qual etapa a especificação é rejeitada se o gatilho não disser o que deve ser feito.\par
(E2) Simule: se $1$ em cada $2$ tarefas caísse na rota (c), qual seria o custo médio adicional por tarefa? E se o Blackboard nunca registrasse voluntários?\par
(E3) Verifique se o seu repositório tem marcação duplicada \cod{ATIVACAO-R613} em algum módulo. Se tiver, o que isso diz sobre a idempotência do gerador?\par
(E4) Liste as $4$ rotas e diga qual delas é acionada quando o usuário digita diretamente o nome de um agente. A ordem de preferência confirma sua resposta?
\end{aviso}

\noindent\textit{Interpretação:} No modelo Blackboard, fontes de conhecimento publicam contribuições em um espaço compartilhado, e um controlador decide qual contribuição altera o estado do problema. Isso separa três operações que uma chamada direta costuma esconder: tornar a tarefa observável, permitir que solucionadores capazes concorram e arbitrar a próxima ação. O exercício testa se o leitor identifica essas fronteiras; a resposta deve ser justificada pelo fluxo descrito no texto, e não apenas pela semelhança entre nomes de funções \citref{NII, H. P. Blackboard systems: the blackboard model of problem solving and the evolution of blackboard architectures. \emph{AI Magazine}, v.~7, n.~2, p.~38-53, 1986}{10.1609/aimag.v7i2.537}{p.~38-53 (p.~42-45, modelo do quadro e ciclo de voluntariado). é o lastro da rota (c): o Blackboard não é um espelho de mensagens, é uma área de trabalho comum onde solucionadores postam, competitorem e o controlador escolhe}{Opportunistic reasoning has the advantage that order of application of the knowledge sources is not predetermined, once it is decided which knowledge sources to activate, the order is determined dynamically}{O raciocínio oportunista tem a vantagem de que a ordem de aplicação das fontes de conhecimento não é pré-determinada; uma vez decidido quais fontes ativar, a ordem é determinada dinamicamente.}{NII, 1986, p.~42--45}
